%%%%%%%%%%%%%%%%%%%%%%%%%%%%%%%%%%%%%%%%%%%%%%%%%%%
%
% File: misc/thesis-commands.tex
%
% Project-specific macros. NOT part of the upstream PSIThesis template.
%
% The thesis runs on a single, non-reusable identifier space
% (docs/01_id_conventions.md). Marking those identifiers semantically
% rather than typing them as literal text buys three things:
%
%   1. one place to change how an identifier looks;
%   2. \index-ability and \label-ability later, without touching prose;
%   3. a grep-able source, so a traceability checker could be pointed at
%      the .tex files as well as at the Markdown.
%
%%%%%%%%%%%%%%%%%%%%%%%%%%%%%%%%%%%%%%%%%%%%%%%%%%%

% ---------------------------------------------------------------------------
% Inline semantic markup
% ---------------------------------------------------------------------------
% These five live in main.tex in the stock template, where they are declared
% after \mainmatter. They are declared here instead so that the front matter
% (abstract, preface, abbreviations) can use them as well.

\newcommand{\keyword}[1]{\textbf{#1}}
\newcommand{\tabhead}[1]{\textbf{#1}}
\newcommand{\code}[1]{\texttt{#1}}
\newcommand{\file}[1]{\texttt{#1}}
\newcommand{\option}[1]{\texttt{\itshape#1}}

% ---------------------------------------------------------------------------
% Identifier space
% ---------------------------------------------------------------------------
% Usage: \hz{01} \sr{003} \cagerule{06} \scn{NOM-01} \met{S1} \dec{69}
%        \gate{4} \phase{F4}
%
% \nbh is a non-breaking hyphen so an identifier never splits over a line
% break -- which matters here because the whole argument of the thesis is
% that these identifiers are atomic.

\newcommand{\nbh}{\mbox{-}}

\newcommand{\hz}[1]{\mbox{H\nbh #1}}
\newcommand{\sr}[1]{\mbox{SR\nbh #1}}
\newcommand{\cagerule}[1]{\mbox{C\nbh #1}}
\newcommand{\scn}[1]{\mbox{SC\nbh #1}}
\newcommand{\met}[1]{\mbox{M\nbh #1}}
\newcommand{\dec}[1]{\mbox{D\nbh #1}}
\newcommand{\gate}[1]{\mbox{G#1}}
\newcommand{\phase}[1]{\mbox{#1}}

% ---------------------------------------------------------------------------
% Verdict labels
% ---------------------------------------------------------------------------
% The campaign aggregator emits three literal verdicts. They are set in the
% monospace face, small caps off, so that a verdict in running prose is
% visibly a machine output and not the author's adjective -- a distinction
% the thesis leans on repeatedly.

\newcommand{\verdictsat}{\textcolor{ubgreen!70!black}{\texttt{SATISFIED}}}
\newcommand{\verdictnot}{\textcolor{ubred}{\texttt{NOT SATISFIED}}}
\newcommand{\verdictind}{\textcolor{gray50}{\texttt{INDETERMINATE}}}
\newcommand{\verdictna}{\textcolor{gray50}{\texttt{N/A}}}
\newcommand{\verdictopen}{\textcolor{gray50}{\texttt{NOT EXECUTED}}}

% Mode names (enforcement / monitoring) -- used often enough to be worth
% a macro, and italicising them consistently keeps the counterfactual legible.
\newcommand{\enf}{\emph{enforcement}}
\newcommand{\mon}{\emph{monitoring}}

% ---------------------------------------------------------------------------
% Evidence-status flags
% ---------------------------------------------------------------------------
% Phase-5 physical evidence is bring-up evidence, never verdict evidence.
% The manuscript says so in prose every time; in this layout the margin is
% the right place for that qualifier, so it never interrupts the sentence.

\newcommand{\prelimflag}{%
  \marginnote{\textcolor{ubred}{\textbf{PRELIMINARY}}\\
  N=1, \mon, unscored. Does not re-score any gate.}}

\newcommand{\retracted}[1]{%
  \marginnote{\textcolor{ubred}{\textbf{RETRACTED}}\\ #1}}

% A short caveat in the margin, without a number. Use for qualifications
% that would otherwise become a parenthesis inside a long sentence.
\newcommand{\caveat}[1]{\marginnote{#1}}

% ---------------------------------------------------------------------------
% Figures
% ---------------------------------------------------------------------------
% The template's \image prepends "../figures/". These wrappers do not, so
% they resolve through \graphicspath and also work for figures/auto/.
%
%   \fig[placement]{width}{file}{caption}{label}
%   \figwide[placement]{file}{caption}{label}   -- spans text + margin column
%   \figmargin[offset]{file}{caption}{label}    -- lives in the margin column
%
% \detokenize on the file argument: every figure in this thesis has an
% underscored name (fig_8_2_safety_invariant). Written straight into the
% document, \includegraphics copes with that; passed through a wrapper macro
% like these it is safer to hand graphicx catcode-12 characters, because the
% underscore has already been tokenised as a subscript by the time the wrapper
% sees it. \detokenize costs nothing when the name is plain.

\newcommand{\fig}[5][t]{%
	\begin{figure}[#1]
		\centering
		\includegraphics[width=#2]{\detokenize{#3}}
		\caption{#4}
		\label{fig:#5}
	\end{figure}
}

\newcommand{\figwide}[4][t]{%
	\begin{figure*}[#1]
		\centering
		\includegraphics[width=\widefigurewidth]{\detokenize{#2}}
		\caption{\label{fig:#4}#3}
	\end{figure*}
}

\newcommand{\figmargin}[4][1]{%
	\begin{marginfigure}[#1\baselineskip]
		\includegraphics[width=\marginparwidth]{\detokenize{#2}}
		\caption{\label{fig:#4}#3}
	\end{marginfigure}
}

% Placeholder for a figure that has not been regenerated on this host.
% Draws a visible frame instead of failing the build, so a missing render
% is obvious in the PDF rather than silent.
%
% Unlike \fig and friends, this macro is the one that TYPESETS the file name
% instead of handing it to \includegraphics, so the name needs protecting:
% every figure here is called fig_something_something, and a raw underscore in
% text mode is "Missing $ inserted". \detokenize re-emits the argument as
% catcode-12 characters, which \texttt sets literally.
%
% \par\medskip rather than \\[...]: inside a fixed-height box a \\ that lands
% next to a \centering can raise "There's no line here to end", and \par is
% unambiguous.
\newcommand{\figmissing}[3][t]{%
	\begin{figure}[#1]
		\centering
		\fbox{%
			\begin{minipage}[c][38mm][c]{0.88\textwidth}
				\centering\sffamily\small
				\textcolor{ubred}{Figure not available on this host}\par\medskip
				\texttt{\detokenize{#2}}\par\medskip
				{\footnotesize regenerate it on the compute host}
			\end{minipage}}
		\caption{#3}
	\end{figure}
}

% ---------------------------------------------------------------------------
% Tables
% ---------------------------------------------------------------------------
% Every table in this thesis is a small evidence table, so they all get the
% same treatment: booktabs rules, tabularx for the wide ones, and a caption
% ABOVE the table (KOMA default), which reads better in a narrow column.

\newcommand{\tabnote}[1]{%
  \par\vspace{0.4em}\noindent{\footnotesize\sffamily #1}}

% The appendix registers -- hazards, requirements, ODD attributes, cage
% parameters, traceability -- are wide tables of sentence-length cells that run
% over a page. They need three things at once that no single stock environment
% gives: page breaking (longtable), the full text+margin width, and a repeating
% header.
%
% longtable cannot live inside table*, and it centres itself in \textwidth by
% default. Setting \LTleft to 0pt and \LTright to \fill pins it flush with the
% text block and lets it run out into the margin column, which is empty on
% these pages. Column widths must sum to \widefigurewidth.
%
% Usage:  \begin{registertable}{@{}L{12mm}L{45mm}...@{}}  ... \end{registertable}
\newenvironment{registertable}[1]
  {\begingroup
   \setlength{\LTleft}{0pt}%
   \setlength{\LTright}{\fill}%
   \scriptsize
   \setlength{\tabcolsep}{3pt}%
   \begin{longtable}{#1}}
  {\end{longtable}\endgroup}

% ---------------------------------------------------------------------------
% \widefigurewidth, made safe for tabularx
% ---------------------------------------------------------------------------
% PSIThesis.cls line 398 defines
%     \def\widefigurewidth{\dimexpr(\marginparwidth + \textwidth + \marginparsep)}
% with NO \relax terminator. That is fine for \includegraphics[width=...],
% where the keyval parser terminates the scan -- which is the only use the
% template itself makes of it.
%
% It is NOT fine as the width argument of tabularx. tabularx measures its
% columns over several passes and does arithmetic on that width repeatedly; an
% unterminated \dimexpr expression is not a length register, and the scan runs
% on into the column specification. Nine tables in this thesis pass
% \widefigurewidth to tabularx.
%
% So: compute it once into a real length and point \widefigurewidth at that.
% Both uses then work, and the class file stays unmodified.
\newlength{\wtw}
\AtBeginDocument{%
  \setlength{\wtw}{\dimexpr\marginparwidth+\textwidth+\marginparsep\relax}%
  \def\widefigurewidth{\wtw}%
}

% A column type that wraps and stays left-aligned inside tabularx.
\newcolumntype{L}[1]{>{\RaggedRight\arraybackslash}p{#1}}
\newcolumntype{Y}{>{\RaggedRight\arraybackslash}X}

% ---------------------------------------------------------------------------
% Units and quantities
% ---------------------------------------------------------------------------
% siunitx is already loaded by the class. These shorthands keep the many
% millimetre / metre-per-second figures consistent.
%
% NOTE on naming: \deg, \m and \sec already exist in LaTeX, and \H is an
% accent. Every shorthand below is therefore prefixed with "q" (quantity)
% so nothing in the kernel or in siunitx is clobbered.

\newcommand{\qmm}[1]{\SI{#1}{\milli\meter}}
\newcommand{\qm}[1]{\SI{#1}{\meter}}
\newcommand{\qmps}[1]{\SI{#1}{\meter\per\second}}
\newcommand{\qdeg}[1]{\SI{#1}{\degree}}
\newcommand{\qhz}[1]{\SI{#1}{\hertz}}
\newcommand{\qs}[1]{\SI{#1}{\second}}
\newcommand{\qpc}[1]{\SI{#1}{\percent}}
